\section{Introduction}
\label{sec:introduction}

Static machine-learning detectors can score Windows portable executable (PE)
files before execution, making them useful for high-volume screening. Public
representations such as EMBER summarize byte distributions, strings, headers,
sections, imports, exports, and data directories in a fixed vector
\cite{anderson2018ember}. Gradient-boosted trees remain strong learners for these
heterogeneous attributes and provide efficient, exact attribution mechanisms
\cite{ke2017lightgbm,lundberg2020local}. Their accuracy nevertheless depends on
the labeled distribution used for training. New software, collection sources,
packers, and malware families can alter that distribution, while expert analysis
of every incoming object is infeasible \cite{ucci2019survey,gibert2020rise,
arp2022dos}.

Pool-based active learning addresses this label bottleneck by selecting a small
batch from an unlabeled pool for an oracle. Uncertainty sampling is a natural
baseline because it selects samples near the current decision boundary
\cite{cohn1996active,settles2009active}. A probability margin, however, collapses
the model's evidence into one scalar. Two samples with the same predicted
probability may differ substantially: one can have uniformly weak contributions,
whereas another can contain large positive and negative effects that nearly
cancel. The latter pattern can reveal a region in which the model combines
conflicting learned associations, but it is invisible to margin-only selection.

We examine whether signed TreeSHAP summaries provide useful acquisition signals.
The proposed family combines predictive uncertainty $U$ with directional entropy
$H$, which measures the balance of positive and negative attribution mass, and
signed-attribution variance $V$, which measures dispersion across feature
coordinates. Five policies---Random, Uncertainty, $U+H$, $U+V$, and a complete
$U+H+V$ policy---are compared with matched initial sets, model configurations,
pools, and label budgets. The evaluation emphasizes the entire learning curve
and labels required to reach a predeclared validation target. This avoids using
only the final model, for which every policy eventually receives the same full
training pool.

Using explanations in acquisition creates a second obligation: the attribution
ranking should be tested against fitted-model behavior rather than accepted from
a visualization. We therefore conduct a separate controlled TreeSHAP evaluation
over 25 fitted models. Sample-specific top-$|\mathrm{SHAP}|$ coordinates are
replaced by statistics computed only from training data and compared with
equal-budget random masks. Mean and median baselines test sensitivity to the
replacement rule, while Spearman and Jaccard agreement measure how feature
rankings change across model seeds and data folds.

The intervention is intentionally described as a feature-space diagnostic. PE
attributes are structured and dependent: replacing isolated histogram, import,
header, or section coordinates need not correspond to a valid executable. A
large deletion effect therefore supports reliance of the fitted model on the
ranked coordinates under a specified intervention. It does not prove that the
features cause malicious behavior, that a binary can be edited accordingly, or
that the explanation is useful to every analyst. We report heterogeneous
diagnostics separately and do not construct an Overall Fidelity Score without a
validated normalization and aggregation rule.

The empirical study uses an identity-audited 2,000-sample subset of EMBER with
2,381 static features, 1,500 benign labels, and 500 malware labels. Each row has
a source index and SHA-256 identity in a locked manifest. The study uses five
outer folds, five active-learning seeds per fold, and five fidelity model seeds
per fold. These controls improve on a single-split pilot, but the data remain a
small, short-period subset evaluated with random row folds. Consequently, the
paper makes no claim about ransomware-specific detection, unseen families,
temporal generalization, or production false-alarm rates.

We study three research questions:

\begin{itemize}
    \item \textbf{RQ1:} Under matched training-label budgets, do acquisition
    policies that combine uncertainty with signed TreeSHAP summaries improve ROC
    and average-precision learning-curve area over random and uncertainty-only
    selection?
    \item \textbf{RQ2:} How do the policies trade learning-curve quality, labels
    required to reach a predeclared validation target, low-FPR behavior, and
    query computation?
    \item \textbf{RQ3:} Under training-mean and training-median replacement, do
    top-$|\mathrm{SHAP}|$ masks affect fitted-model ranking more than equal-count
    random masks, and how stable are the rankings across folds and model seeds?
\end{itemize}

The contributions are:

\begin{itemize}
    \item An explanation-guided acquisition family that combines prediction
    margin with directional balance and dispersion of signed TreeSHAP
    contributions, together with component ablations under matched budgets.
    \item A multi-fold, multi-seed empirical protocol that reports ROC/AP area
    under the learning curve (AULC), labels to a validation-defined target,
    low-FPR summaries, paired fold effects, and computation. The results expose a
    useful trade-off: $U+H$ has the strongest mean AULC, $U+V$ reaches the target
    with the fewest labels, and no variant wins every criterion.
    \item A repeated controlled evaluation of TreeSHAP rankings using two
    training-derived baselines, equal-budget random controls, three deletion
    budgets, 25 fitted models, and within-/between-fold stability analysis.
    \item A claim-bounded reproducibility design that links dataset identity,
    splits, model seeds, query decisions, predictions, masks, and reported replay
    checks while separating computational reproducibility from external
    generalization.
\end{itemize}

Section~\ref{sec:related} reviews static PE learning, batch active learning, and
attribution evaluation. Section~\ref{sec:problem} defines the estimands and claim
boundaries. Sections~\ref{sec:method} and~\ref{sec:protocol} specify the methods
and protocol. Sections~\ref{sec:results}--\ref{sec:threats} present results,
interpretation, and validity threats; Section~\ref{sec:conclusion} concludes.

